\section*{\S\,2. Диаграммы Эйлера}
\addcontentsline{toc}{section}{\S\,2. Диаграммы Эйлера}
\markright{\S\,2. Диаграммы Эйлера}
\subsection*{Условия}
1. Используя диаграммы Эйлера, доказать тождества\fixed{в оригинале последнее тождество напечатано как
$A\cup[(B\setminus\ov A)\sd C]\equiv(A\cap B)\sd C$, которое тождеством
не является (при $A=I$, $B=C=\emptyset$ левая часть равна $I$, правая "---
$\emptyset$). Восстановленный вариант и разбор "--- в решении ниже.}
\[
\begin{gathered}
(A\cup C)\cap B\equiv(A\cap B)\cup(\ov A\cap B\cap C);\\
[A\cap(B\cup C)]\cup(\ov A\cap C)\equiv C\cup(A\cap B);\\
(A\cap C)\cup(B\cap C)\cup(C\setminus A)\equiv C;\qquad A\sd B\equiv(A\cup B)\setminus(A\cap B);\\
[A\cap(B\sd C)]\cup(B\sd C)\equiv B\sd C;\\
A\cap[(B\setminus\ov A)\sd C]\equiv(A\cap B)\sd(A\cap C).
\end{gathered}
\]


2. Выпуклые области $A,B,C\subset\R^2$ разбивают плоскость на восемь частей,
различающихся состояниями принадлежности их точек множествам $A$, $B$ и $C$.
Может ли добавление выпуклой области $D$ привести к разбиению $\R^2$ на 16
подобластей указанного типа? Какое максимальное число таких подобластей
могут образовать $n$ выпуклых областей в $\R^2$?

3. Существуют ли (для всякого $n$) конгруэнтные (т.\,е. получающиеся одна из
другой переносом и поворотом) области $A_1,\dots,A_n\subset\R^2$,
разбивающие плоскость на максимально возможное ($2^n$) число частей, которые
различаются между собой состояниями принадлежности их точек областям
$A_1,\dots,A_n$. (Заметим, что для $n\le3$ такими областями, например,
являются круги одинакового радиуса.)

\subsection*{Решения}
\begin{solution}[1 (\S\,2)]
Нумеруем конституентные области трех кругов, как на рис.~3,~\emph{б},
и записываем, какие из них образуют левую и правую части. Удобнее
обозначать область тройкой $\alpha\beta\gamma$ состояний принадлежности
к $A$, $B$, $C$ ($1$ "--- внутри, $0$ "--- вне); тогда
\[
1=111,\ 2=110,\ 3=011,\ 4=101,\ 5=010,\ 6=001,\ 7=100,\ 8=000 .
\]
\begin{center}\small
\begin{tabular}{lll}
\toprule
тождество & левая часть & правая часть\\
\midrule
$(A\cup C)\cap B\equiv(A\cap B)\cup(\ov A\cap B\cap C)$ & $\{1,2,3\}$ & $\{1,2\}\cup\{3\}$\\
$[A\cap(B\cup C)]\cup(\ov A\cap C)\equiv C\cup(A\cap B)$ & $\{1,2,4\}\cup\{3,6\}$ & $\{1,3,4,6\}\cup\{1,2\}$\\
$(A\cap C)\cup(B\cap C)\cup(C\setminus A)\equiv C$ & $\{1,4\}\cup\{1,3\}\cup\{3,6\}$ & $\{1,3,4,6\}$\\
$A\sd B\equiv(A\cup B)\setminus(A\cap B)$ & $\{4,7,3,5\}$ & $\{1,\dots,5,7\}\setminus\{1,2\}$\\
$[A\cap(B\sd C)]\cup(B\sd C)\equiv B\sd C$ & $\{2,4\}\cup\{2,4,5,6\}$ & $\{2,4,5,6\}$\\
\bottomrule
\end{tabular}
\end{center}
Во всех строках множества номеров совпадают, а так как три круга
в положении рис.~3,~\emph{б} реализуют все восемь логически возможных
состояний, совпадение областей на диаграмме и означает тождество
(см. также упр.~2 \S\,3 и решение упр.~7 \S\,3, где это рассуждение
обосновано строго). Четвертое и пятое тождества видны уже на диаграмме
двух кругов (пятое "--- это поглощение $(X\cap Y)\cup Y\equiv Y$).

\emph{Последнее тождество.} Поскольку $B\setminus\ov A\equiv B\cap A$,
напечатанное в книге равенство
$A\cup[(B\setminus\ov A)\sd C]\equiv(A\cap B)\sd C$ имеет левую часть
$A\cup\{1,2\}\sd\{1,3,4,6\}=\{1,2,4,7\}\cup\{2,3,4,6\}=\{1,2,3,4,6,7\}$,
а правую $\{2,3,4,6\}$: они различаются областями $1$ и $7$ (например, при
$A=I$, $B=C=\emptyset$). Значит, в книге опечатка. Естественный
верный вариант "--- с пересечением вместо объединения:
\[
A\cap[(B\setminus\ov A)\sd C]\equiv(A\cap B)\sd(A\cap C):
\]
слева $\{1,2,4,7\}\cap\{2,3,4,6\}=\{2,4\}$, справа
$\{1,2\}\sd\{1,4\}=\{2,4\}$. Это тождество выражает
дистрибутивность $\cap$ относительно $\sd$:
$A\cap(X\sd Y)\equiv(A\cap X)\sd(A\cap Y)$. (Верно также
$A\cap[(B\setminus\ov A)\sd C]\equiv(A\cap B)\sd(A\cap C)\equiv A\cap(B\sd C)$.)
\end{solution}

\begin{solution}[2 (\S\,2)]
\emph{Ответ: да; $n$ выпуклых областей могут дать все $2^n$
логически возможных частей.}

Для $n=4$ это показывает, например, известная диаграмма Венна из четырех
равных эллипсов (эллипс выпуклый). Общая конструкция для любого $n$:
возьмем на окружности $2^n$ различных точек $q_\alpha$, занумерованных
наборами $\alpha=(\alpha_1,\dots,\alpha_n)\in\{0,1\}^n$, и положим
\[
A_i=\operatorname{conv}\{q_\alpha:\ \alpha_i=1\}
\]
(выпуклая оболочка; чтобы получить область с внутренними точками, ее можно
слегка <<раздуть>>, взяв $\varepsilon$-окрестность при малом
$\varepsilon$ "--- выпуклость сохраняется). Точки на окружности находятся
в выпуклом положении: ни одна из них не лежит в выпуклой оболочке
остальных и даже отстоит от нее на положительное расстояние. Поэтому
$q_\alpha\in A_i\iff\alpha_i=1$, т.\,е. точка $q_\alpha$ принадлежит
конституенте $A_1^{\alpha_1}\cap\dots\cap A_n^{\alpha_n}$, и все $2^n$
конституент непусты (а вместе с малыми окрестностями точек $q_\alpha$,
лежащими вне всех $A_i$ с $\alpha_i=0$, "--- имеют внутренние точки).

Итак, максимальное число частей, различающихся состояниями
принадлежности, равно $2^n$ "--- больше быть не может, так как наборов
состояний всего $2^n$. (Связными эти части, вообще говоря, не будут; если
же считать связные компоненты, то их число для выпуклых областей не
ограничено никакой функцией от $n$: границы двух выпуклых многоугольников
могут пересекаться в сколь угодно большом числе точек.)
\end{solution}

\begin{solution}[3 (\S\,2)]
\emph{Ответ: да, для всякого $n$; можно даже обойтись одними
параллельными переносами.}

Выберем $n$ векторов $v_1,\dots,v_n$ и $2^n$ точек $x_\alpha$,
$\alpha\in\{0,1\}^n$, <<в общем положении>>, так чтобы все $n2^n$ точек
$x_\alpha-v_i$ были попарно различны (этого можно добиться сколь угодно
малым шевелением). Разобьем их на два конечных множества
\[
P_+=\{x_\alpha-v_i:\ \alpha_i=1\},\qquad P_-=\{x_\alpha-v_i:\ \alpha_i=0\}.
\]
Пусть $S$ "--- связная открытая область, содержащая все точки $P_+$ и не
содержащая ни одной точки $P_-$ (например, объединение малых кругов вокруг
точек $P_+$, соединенных тонкими <<коридорами>>, обходящими конечное
множество $P_-$, "--- плоскость без конечного множества точек связна).
Положим $A_i=S+v_i$. Тогда
\[
x_\alpha\in A_i\iff x_\alpha-v_i\in S\iff\alpha_i=1,
\]
т.\,е. $x_\alpha$ лежит в конституенте с номером $\alpha$, и все $2^n$
конституент непусты (и открыты в окрестности $x_\alpha$, если круги
выбрать малыми). Области $A_i$ конгруэнтны "--- они получаются друг из
друга переносом.

Для $n\le3$ подходят равные круги (рис.~3,~\emph{б}). Для $n\ge4$
равные (и вообще любые) \emph{круги} уже не годятся: $n$ окружностей
разбивают плоскость не более чем на $n^2-n+2$ связных частей (каждая новая,
$k$-я, окружность пересекает прежние не более чем в $2(k-1)$ точках
и добавляет не более $2(k-1)$ частей). Каждая непустая конституента
содержит хотя бы одну такую часть, поэтому непустых конституент не больше
$n^2-n+2$, а $n^2-n+2<2^n$ при $n\ge4$. Поэтому при $n\ge4$ приходится
брать области более сложной формы, как в построении выше. Для простых $n$
существуют даже симметричные диаграммы Венна из $n$ конгруэнтных кривых,
переходящих друг в друга поворотом (Григгс, Киллиан, Сэвидж, 2004).
\end{solution}
